\documentclass[12pt,a4paper]{article}

\usepackage[utf8]{inputenc}
\usepackage[T1]{fontenc}
\usepackage[margin=2.5cm]{geometry}
\usepackage{hyperref}
\usepackage{microtype}
\usepackage{parskip}
\usepackage{lmodern}

\hypersetup{colorlinks=true, urlcolor=blue!60!black}

\pagestyle{empty}

\begin{document}

\begin{flushright}
Kundai Farai Sachikonye \\
Auf der Lichtung 19, 82178 Puchheim \\
\href{https://github.com/fullscreen-triangle}{github.com/fullscreen-triangle} \\
\texttt{kundai.sachikonye@bitspark.com} \\
(+49) 1626386344 \\[6pt]
16 May 2026
\end{flushright}

\vspace{1em}

Leibniz-Institut für Analytische Wissenschaften --- ISAS --- e.V. \\
Arbeitsgruppe Mehrdimensionale OMICS-Datenanalyse \\
Dortmund, Germany \\
Reference: 398\_2026

\vspace{1.5em}

\textbf{Re: Doktorand:in (m/w/d) --- Mehrdimensionale OMICS-Datenanalyse}

\vspace{1em}

Dear Hiring Committee,

I am applying for the doctoral position in the Multi-Dimensional OMICS Data Analysis
research group. The combination of modalities in this position --- phenotypic,
microscopic, PET/MRI, MALDI MSI, and LC-MS/MS data --- is precisely what I have
built an integrated computational framework to handle. The relevant tools are publicly
deployed and available without installation:

\begin{itemize}
  \item Genomics and sequence analysis:
    \href{https://vivid-symbolism.vercel.app/search}{vivid-symbolism.vercel.app/search}
  \item Mass spectrometry (MALDI MSI / LC-MS) experiment design and simulation:
    \href{https://lavoisier.vercel.app/experiment}{lavoisier.vercel.app/experiment}
  \item Microscopy image analysis (PNG / JPG / TIFF, GPU-accelerated):
    \href{https://hieronymus-omega.vercel.app/observe/microscopy}{hieronymus-omega.vercel.app/observe/microscopy}
  \item Protein analysis (folding, trajectory, catalysis, dynamics):
    \href{https://shakespear-nine.vercel.app/instrument}{shakespear-nine.vercel.app/instrument}
  \item Cheminformatics (molecular spectroscopy, docking, fingerprints, 3D structure, GPU):
    \href{https://honjo-masamune.vercel.app/tools}{honjo-masamune.vercel.app/tools}
\end{itemize}

\textbf{Multi-dimensional OMICS integration.}
The \href{https://github.com/fullscreen-triangle/gospel}{Gospel} framework processes
whole-genome variant data at $10^6$ variants per second, integrates RNA-seq
quantification, and applies Bayesian network inference and Gaussian mixture models
for multi-omics statistical modelling. It has been validated against 1000 Genomes
Phase~3, gnomAD v3.1.2, and UK Biobank, and covers the full pipeline from
variant calling through pathway analysis (KEGG, Reactome) to functional annotation.
The \href{https://github.com/fullscreen-triangle/lavoisier}{Lavoisier} framework
provides the MALDI MSI and LC-MS/MS layer: a dual pipeline combining numerical
MS1/MS2 analysis with a CNN-based visual pipeline (PyTorch) that converts spectra
into temporal image sequences, achieving 98.9\,\% feature extraction accuracy
against the NIST reference library and supporting memory-mapped I/O for datasets
exceeding 100\,GB. For the microscopy modality, the
\href{https://github.com/fullscreen-triangle/helicopter}{Helicopter} platform
fuses fluorescence, brightfield, phase-contrast, and spectral channels through
sequential constraint exclusion across twelve measurement modalities, reducing
structural ambiguity from ${\sim}10^{60}$ configurations to a unique assignment.
Together these five systems address all major data types listed in this position
within a shared S-entropy coordinate framework that makes cross-modal registration
tractable. The
\href{https://github.com/fullscreen-triangle/shakespear}{Shakespear} framework
adds protein-level analysis --- folding, trajectory, catalysis, and dynamics via
Kuramoto oscillator synchronization in S-entropy space --- directly relevant to
proteomics readouts and protein-state transitions in the multi-omics pipeline.
The \href{https://honjo-masamune.vercel.app/tools}{Honjo-Masamune} platform
provides zero-parameter GPU-accelerated molecular spectroscopy (vibrational modes,
docking, 3D structure, fingerprints) running entirely in the browser, addressing
the metabolite and small-molecule identification layer.

\textbf{AI-based personalized risk assessment.}
This is where the theoretical work I have developed is most directly applicable.
The \href{https://github.com/fullscreen-triangle/syndrome}{Syndrome} framework
computes personalised disease state as a vector $\mathbf{D}\in[0,1]^8$ across
eight cellular oscillator classes (protein, enzyme, channel, membrane, ATP, genetic
expression, calcium, circadian), each with a coherence index $\eta_i\in[0,1]$.
The framework locates the patient's current state within a partition geometry on
the $[0,1]^8$ hypercube, computes an $\mathcal{O}(N\log N)$ trajectory to the
healthy attractor, and selects the minimum-side-effect therapeutic intervention
via sparse $\ell_1$ linear programming. This is a personalized risk assessment
model in the sense the position describes: patient-specific state, validated
biomarker readouts as coherence indices, and a computed trajectory rather than
a population-level risk score.

A companion paper, \textit{On Disease Epistemology in Blind Receiver} (affiliated
with the AIMe Registry and TUM School of Life Sciences Weihenstephan), proves that
loop-holonomy self-consistency is the unique viable disease diagnostic criterion,
achieving AUC$\,=1.00$ versus 0.896 for template comparison with 25/25 validation
experiments passing, and providing a sparse LP for therapeutic target selection.

\textbf{Large language models and AI tooling.}
The position specifically lists LLM experience. Beyond using LLM APIs, I have
formalised a training principle for domain-specific models:
the \href{https://github.com/fullscreen-triangle/purpose}{Purpose} framework proves
that training exclusively at the hardest single-objective regime (the pure-intent
regime) yields competence across all constrained sub-regimes by feasibility inclusion
rather than generalisation --- with $(N{+}1)\times$ sample savings and an exponential
$2^{N-1}\times$ saving for binary-nested constraint cascades.
The \href{https://github.com/fullscreen-triangle/four-sided-triangle}{Four-Sided-Triangle}
system implements an 8-stage RAG pipeline with adaptive hybrid LLM and mathematical
solver selection and a Rust backend providing 10--50$\times$ speedup for
compute-intensive stages. For the OMICS context this means: not deploying a
general-purpose model on specialised data, but training or selecting a model whose
capability envelope was defined by the hardest instances of the target domain.

\textbf{Background.}
I hold an MSc in Computational Biology (Universität Tübingen) and an MSc in
Pharmaceutical Biotechnology (HAW Hamburg), and I completed doctoral research in
Computational Lipidomics at TUM / Bayerisches Forschungsinstitut, covering mass
spectrometry pipelines, federated learning for distributed clinical analysis, and
multi-omics modelling. I am legally resident in Germany with full work authorisation
and am available immediately. English is native; German is conversational.

\vspace{1.5em}
Yours sincerely,

\vspace{2.5em}
\textbf{Kundai Farai Sachikonye}

\end{document}
